\section{Discussion}
\label{sec:discussion}

\paragraph{What kind of thing is depth structure?}
The composite picture is specific enough to be useful. Residual updates are
predictable across depth, but the predictability decomposes into (i) a
share carried by slow drift along a few persistent state directions, which
we set aside as uninformative, (ii) a genuine
increment-level structure that a randomly initialised network already
exhibits --- substantially so on \pythia{}, where the untrained network in
fact exceeds the trained one; about half of it on GPT-2-small ---
and (iii) a learned component that only becomes visible in the
\emph{generative} fit --- roughly half the SSM's reconstruction advantage,
majority at sparse operating points. Meanwhile the dynamics are pervasively
non-stationary in depth, first-order, and --- per the fitted model's
gauge-invariant state path --- attenuating by half per layer of gap, a
one-layer memory (measured with sources at every layer, which leaves the
state little to carry; \cref{sec:evaluation:operator}). A transformer, on
this evidence, does
not implement a depth-shared propagation rule; it implements a different
small-memory rule at every layer, over features that persist via the additive
stream itself rather than via the incremental dynamics. The Granger reach of
20 layers coexisting with $\kstar\!\approx\!1$ stops being a paradox once
split by sublayer: a small, deep-reaching attention channel superposed on an
energy-dominant, memoryless MLP channel. (That split localises the reach; it
does not by itself certify the AR structure as persistence-free, since both
sublayer writes read the same drifting state --- \cref{sec:phase1:sublayer}.)

\paragraph{Consequences for cross-layer interpretability methods.}
H2's failure supports CLT's design decision not to share decoders along
depth \citep{ameisen2025circuit} ---
our attempt to compress that family along depth was the pre-registered
falsifier, and it failed in CLT's favour. That is support for
depth-indexing, which a first-order state has too. The case for a free
decoder at every (source,\,target) pair comes from the quantitative
comparison, and it is budget-conditional: the free decoders buy a real
margin on every model we tested. On the primary model that margin is small
and data-bought, and a first-order state recovers about four fifths of the
CLT's gain over a rank-matched per-layer decoder at a sixth of the
parameters and generalises better once
the training set falls below $10^3$ examples. On the smaller replication
models the state recovers less of the CLT's gain, as little as 30\% on
GPT-2-small (\cref{sec:evaluation:clt}). A per-layer decoder widened to the
state's parameter count recovers as much as the state does, so on frozen
sources the parameter saving comes from not fitting a decoder per layer
pair, not from the state. For practitioners the operational
reading is therefore conditional: where the free decoders cost several times
the parameters, as on our primary model, they are the poor spend under a
parameter budget, and under a data budget the depth-indexed state beat them
(we did not run the wide per-layer decoder there); if you need maximum
fidelity and
have data, free per-pair decoders are not redundant. Our stability finding
adds an orthogonal
caution: the SSM's write directions for individual features were not
reproduced across refits over TopK dictionaries in any configuration we
tested beyond the six-layer debug lane, while over L1
dictionaries they were reproduced everywhere (perfectly so on the primary
model). What we re-initialised is the downstream model, not the SAE, so the
caution we can support is about downstream fits: a pipeline that learns
per-feature readouts over a TopK dictionary and then names, steers, or
ablates \emph{individual} features should check those readouts survive a
refit. Whether the dictionaries themselves survive retraining is a separate
question \citep{paulo2026seeds,song2026consistency}; in our small-data
regime they do not, under either mechanism (\cref{sec:evaluation:stability}),
so the downstream caution adds to the dictionary-level one rather than
replacing it.

\paragraph{Consequences for factual-recall stories.}
The causal null is not a failure to find structure; it is a measurement of
what kind of structure recall has. Single features move the answer by at most
an eighth of the pre-registered threshold under the TopK dictionary, and by
about half of it under the L1 dictionary; the effect that does clear
threshold
is reproduced by a norm-matched random direction; and the larger
single-feature effects under L1 still do not localise it. This is a
controlled statement, the same under both dictionaries we tried, that recall
at the query position is a
distributed, magnitude-driven code --- consistent with the drift in the
editing literature away from single-component accounts
\citep{hase2023does}, and a calibration point for ablation-based evaluation
generally: group-ablation effects need magnitude-matched \emph{direction}
controls before they can be read as evidence about privileged feature sets
--- the intervention-level analogue of the random-baseline sanity checks
\citet{korznikov2026sanity} argue SAE evaluations need.

\paragraph{Consequences for probing.}
The state-probe result cuts the other way: trained purely to reconstruct
updates, the 32-dimensional state retains answer-identity information that
variance-optimal compressions of the same codes discard. (The comparison is
among summaries of one set of codes. Those codes are fit to each example's
whole trajectory, and what they pick up about later updates enters through
the state map, so the comparison says nothing about the state against the
raw stream and may flatter it against the code summaries;
\cref{sec:evaluation:probe}.) If the lead survives codes inferred from
earlier layers only, it says that reconstruction
objectives, at least here, learn compressions aligned with what the circuit
is computing --- a modest positive datum for the program of
extracting functional summaries from decompositions, and one that would pair
naturally with the fast-attenuating state path: a short-memory,
information-dense
running summary is exactly what a fast-forgetting state can be.

\paragraph{The EMG intuition, in retrospect.}
The convolutive-source-separation analogy \citep{holobar2007multichannel}
supplied H2, and H2 is false: transformer depth is not a homogeneous
propagation medium. What survives from the analogy is the discipline it
imported --- model the increments, put sources at every position, and treat
the mixing structure as an estimable, testable object. The estimated object
just turned out to be depth-indexed, fast-attenuating, and circuit-specific
rather than shared.

\paragraph{Future work.}
Three items are specified and costed in the project record: deconfounding the
depth series from width (Pythia-1b/1.4b/2.8b with one dictionary type, three
seeds); the two missing families (Gemma-2's sandwich normalisation,
Llama-3.1-8B) which also enable the pre-registered $\ge$1{,}000-per-class
relation test; and a natively-trained CLT. Two more suggest themselves from
the results: whether the state-path decay rate and the mid-depth
functional-state window are universal across scale, and whether the sublayer
AR structure survives the T2 surrogate and the untrained control that
\cref{sec:phase1:sublayer} did not run. Four checks that this paper names and
does not run fit the study GPU: the state probe with codes inferred from
earlier layers only and with a feedthrough-image arm
(\cref{sec:evaluation:probe}); the source-prefix
sweep with its bias initialised at the per-depth mean
(\cref{sec:evaluation:prefix}); the parameter-matched comparison at further
seeds, over the L1 dictionary, at the other depths and on the untrained
network (\cref{sec:evaluation:gate3}); and the transfer of AR kernels between
the trained and the untrained network (\cref{sec:phase1:pressure}).
